\documentclass[11pt]{article}

\usepackage[T1]{fontenc}
\usepackage[a4paper,margin=2.7cm]{geometry}
\usepackage{amsmath,amssymb,amsthm,mathtools}
\usepackage{xcolor}
\usepackage{hyperref}
\usepackage{microtype}

\hypersetup{colorlinks=true,linkcolor=blue!55!black,citecolor=blue!55!black,
  urlcolor=blue!55!black}

\theoremstyle{plain}
\newtheorem{theorem}{Theorem}[section]
\newtheorem{lemma}[theorem]{Lemma}
\newtheorem{proposition}[theorem]{Proposition}
\newtheorem{corollary}[theorem]{Corollary}
\theoremstyle{definition}
\newtheorem{definition}[theorem]{Definition}
\theoremstyle{remark}
\newtheorem*{remark}{Remark}
\newtheorem*{status}{Status}

\newcommand{\C}{\mathbb{C}}
\newcommand{\R}{\mathbb{R}}
\newcommand{\Q}{\mathbb{Q}}
\newcommand{\Z}{\mathbb{Z}}
\newcommand{\N}{\mathbb{N}}
\newcommand{\Alg}{\overline{\Q}}
\newcommand{\Tc}{\mathcal{T}}
\newcommand{\PAlg}{\widehat{\Alg}}
\newcommand{\PGL}{\operatorname{PGL}}
\newcommand{\J}{\mathrm{J}}
\newcommand{\lean}[1]{\textup{\texttt{#1}}}

\title{The conjugation degree on the transcendental locus\\
\large from the ground up}
\author{Carlo Perassi}
\date{}

\begin{document}
\maketitle

\begin{abstract}
A ground-up account of the conjugation degree
$\delta(z)=[\Alg(z,\bar z):\Alg(z)]$ on $\Tc=\C\setminus\Alg$, written to be
formalized. Every statement below is marked with its status in the Lean
development at \url{https://github.com/carlok/moebius-transcendental-lean}:
either the name of the declaration that proves it, or an explicit note that it
is stated and not proved. No claim is made for any statement not so marked.
\end{abstract}

\section{Setting}

Write $\Alg\subset\C$ for the field of algebraic numbers and
\[
  \Tc:=\C\setminus\Alg
\]
for the transcendental locus. Let $\PAlg=\Alg\cup\{\infty\}$.

\begin{status}
$\Alg$ is \lean{Qbar}, defined as
\lean{(algebraicClosure Q C).toSubfield}, with \lean{mem\_Qbar\_iff}
identifying membership with \lean{IsAlgebraic Q}. $\Tc$ is
\lean{Transcendentals}. Both proved.
\end{status}

Let $\Gamma=\PGL(2,\Alg)$ act on $\widehat\C$ by
$z\mapsto(az+b)/(cz+d)$ with $a,b,c,d\in\Alg$, $ad-bc\ne0$, and let
$\J(z)=\bar z$. Put $\Gamma^{\pm}=\langle\Gamma,\J\rangle$; its elements
outside $\Gamma$ are $z\mapsto(a\bar z+b)/(c\bar z+d)$ with algebraic
coefficients.

\section{The invariant}

\begin{definition}\label{def:delta}
For $z\in\C$, the \emph{conjugation degree} is
\[
  \delta(z):=[\Alg(z,\bar z):\Alg(z)]\in\N\cup\{\infty\},
\]
the degree of $\bar z$ over the field generated by $\Alg$ and $z$.
\end{definition}

\begin{status}
\lean{conjDegree}, valued in \lean{ENat}, that is $\N\cup\{\infty\}$, defined as the degree of the minimal
polynomial of $\bar z$ over $\Alg(z)$ when that element is integral, and $\top$
otherwise. The field $\Alg(z)$ is \lean{gen}. The strata
$\Tc_n=\{z\in\Tc:\delta(z)=n\}$ are \lean{stratum}. Definitions only; the
results of \S3 are stated but not yet proved.
\end{status}

\begin{remark}
The definition is by minimal polynomial rather than by
\lean{Module.finrank}. Two reasons: it avoids the intermediate-field tower
instances, and $\top$ is a genuine value where \lean{finrank} would return the
junk value $0$ for an infinite degree, silently merging $\Tc_\infty$ with the
empty case.
\end{remark}

\section{Elementary consequences}

Everything in this section is elementary; it is recorded because the
formalization needs it, not because it is difficult.

\begin{lemma}[No holomorphic stabilizers]\label{lem:stab}
If $g\in\Gamma$ is not the identity, every fixed point of $g$ lies in $\PAlg$.
\end{lemma}

\begin{proof}
Write $g(z)=(az+b)/(cz+d)$ with $a,b,c,d\in\Alg$. A finite fixed point
satisfies $cz^2+(d-a)z-b=0$, a polynomial with coefficients in $\Alg$, not
identically zero since $g\ne\mathrm{id}$; so $z\in\Alg$. The remaining
projective point is $\infty\in\PAlg$.
\end{proof}

\begin{proposition}[Transcendence is preserved]\label{prop:preserve}
If $z\in\Tc$ and $g\in\Gamma^{\pm}$ then $g(z)\in\Tc$.
\end{proposition}

\begin{proof}
Suppose $g\in\Gamma$ and $g(z)\in\Alg$. Then $g^{-1}\in\Gamma$ has algebraic
coefficients and $z=g^{-1}(g(z))\in\Alg$, contradicting $z\in\Tc$. For
$g\notin\Gamma$ write $g=h\circ\J$ with $h\in\Gamma$; complex conjugation
preserves $\Alg$, hence preserves $\Tc$, and the holomorphic case applies
to $h$.
\end{proof}

\begin{proposition}[$\delta$ is $\Gamma^{\pm}$-invariant]\label{prop:inv}
$\delta(g(z))=\delta(z)$ for every $z\in\Tc$ and $g\in\Gamma^{\pm}$.
\end{proposition}

\begin{proof}
Let $g\in\Gamma$ with coefficients $a,b,c,d\in\Alg$. Since $g$ and $g^{-1}$
have algebraic coefficients, $\Alg(g(z))=\Alg(z)$. Conjugating the defining
expression gives $\overline{g(z)}=\bar g(\bar z)$, where $\bar g$ has the
conjugated coefficients $\bar a,\bar b,\bar c,\bar d$, again in $\Alg$ because
$\Alg$ is conjugation-stable. Hence
$\Alg(g(z),\overline{g(z)})=\Alg(z,\bar z)$, and the two degrees agree. For
$g=\J$ the two generators are exchanged and the field is unchanged.
\end{proof}

\begin{proposition}[Degree one]\label{prop:deg1}
For $z\in\Tc$, $\delta(z)=1$ if and only if $\bar z\in\Alg(z)$.
\end{proposition}

\begin{proof}
Immediate from Definition~\ref{def:delta}: a degree-one extension is trivial.
\end{proof}

\begin{remark}[Why degree one is the interesting stratum]
If $z\bar z=\rho$ with $\rho\in\Alg$ and $z\ne0$, then $\bar z=\rho/z$ lies in
$\Alg(z)$, so $\delta(z)=1$. Complex conjugation on $\Alg(z)$ is then
determined by the field operations rather than being independent data. That
case is formalized separately, machine-checked and kernel-verified, at
\url{https://github.com/carlok/diaz-modulus-lean}.
\end{remark}

\section{Multiplicity on a plane curve}

For the target of \S5 one needs the multiplicity of a point on a plane curve.
Mathlib has no such notion; it is assembled here from
\lean{MvPolynomial.homogeneousComponent}.

\begin{definition}\label{def:mult}
For $P\in K[X,Y]$ and $p\in K^2$, let $P_p$ denote $P$ translated so that $p$
is the origin. The \emph{multiplicity} of $p$ on $P=0$ is
\[
  \operatorname{mult}_p(P):=\min\{\,n:\text{the degree-$n$ homogeneous
  component of }P_p\text{ is non-zero}\,\}.
\]
\end{definition}

\begin{status}
\lean{shiftTo} and \lean{mult}. Proved.
\end{status}

\begin{lemma}\label{lem:mult-wd}
If $P\ne0$ then the set in Definition~\ref{def:mult} is non-empty, so the
minimum exists; and $P_p\ne0$.
\end{lemma}

\begin{proof}
A polynomial is the sum of its homogeneous components, so if all of them
vanish then so does the polynomial; that gives the first claim. Translation by
$p$ is inverted by translation by $-p$, so it is injective, giving the second.
\end{proof}

\begin{status}
\lean{nonempty\_of\_ne\_zero} and \lean{shiftTo\_ne\_zero}. Proved,
\lean{sorry}-free.
\end{status}

\section{The target, stated and not proved}

\begin{theorem}[Conjugation degree is circular degree]\label{thm:target}
Let $V$ be an absolutely irreducible plane curve over $\Alg\cap\R$ whose real
locus contains a transcendental point, let $\overline V\subset\mathbb P^2_{\Alg}$
be its projective closure of degree $d$, and let
$m=\operatorname{mult}_{I_+}(\overline V)=\operatorname{mult}_{I_-}(\overline V)$
at the circular points $I_\pm=[1:\pm i:0]$, zero when the point is absent.
Then every $z=x+iy\in V(\R)\cap\Tc$ satisfies
\[
  \delta(z)=d-m.
\]
\end{theorem}

\begin{status}
\textbf{Not proved, and not currently provable in Lean.} The argument requires
B\'ezout's theorem for plane curves, the normalization of a projective curve,
and the identification of the degree of a morphism with a function-field
extension degree. As of Mathlib \texttt{81a5d257}, none of the three is
available: the \texttt{Bezout} declarations in Mathlib are \texttt{IsBezout},
the ring-theoretic class asserting that finitely generated ideals are
principal, and there is no \texttt{intersectionMultiplicity}. The statement is
recorded here so the formalization has a target and so the gap is visible.
\end{status}

\begin{remark}[A shortcut that does not work]
It is tempting to read $m$ off the leading form of $P$, as the multiplicity of
$x^2+y^2$ as a factor of the degree-$d$ part. That quantity is the
intersection multiplicity of $\overline V$ with the line at infinity, which
bounds $\operatorname{mult}_{I_+}(\overline V)$ from above and equals it only
when the branch is transverse. It agrees in the cases one checks by hand — a
circle has leading form $x^2+y^2=(y-ix)(y+ix)$, so $d=2$, $m=1$; a real line
has $d=1$, $m=0$ — because both are smooth at infinity. It is not a proof.
\end{remark}

\section{Status}

\begin{center}
\begin{tabular}{llll}
\hline
Statement & Lean declaration & Module & Status\\
\hline
$\Alg$ & \lean{Qbar}, \lean{mem\_Qbar\_iff} & \lean{Basic} & proved\\
$\Tc$ & \lean{Transcendentals} & \lean{Basic} & proved\\
$\Alg(z)$ & \lean{gen} & \lean{Basic} & proved\\
Def.~\ref{def:delta} & \lean{conjDegree} & \lean{Basic} & defined\\
$\Tc_n$ & \lean{stratum} & \lean{Basic} & defined\\
Lem.~\ref{lem:stab} & --- & --- & stated only\\
Prop.~\ref{prop:preserve} & --- & --- & stated only\\
Prop.~\ref{prop:inv} & --- & --- & stated only\\
Prop.~\ref{prop:deg1} & --- & --- & stated only\\
Def.~\ref{def:mult} & \lean{shiftTo}, \lean{mult} & \lean{Multiplicity} & proved\\
Lem.~\ref{lem:mult-wd} & \lean{nonempty\_of\_ne\_zero}, \lean{shiftTo\_ne\_zero} & \lean{Multiplicity} & proved\\
Thm.~\ref{thm:target} & --- & --- & not provable yet\\
\hline
\end{tabular}
\end{center}

Five declarations are proved and \lean{sorry}-free. Two are definitions with no
results yet attached. Four elementary statements are written down but not
formalized. One theorem is out of reach and the reason is named.

\section*{Use of AI assistance}

The mathematics, this note and the Lean development were produced with
substantial assistance from large language models, directed and accepted by the
author, who is responsible for the content. Every declaration marked proved
above is machine-checked by Lean~4 at the pinned toolchain
\texttt{leanprover/lean4:v4.32.0}; no statement marked proved rests on an
assertion by a model, and no statement not so marked is claimed.

\end{document}
